\section{Konzeption der Artefakte}
\label{sec:Konzeption}

% Ziel: 6-8 Seiten
% Inhalt: Fragebogen, Mapping-Workbook, SOPs

\subsection{GRACE-Onboarding-Fragebogen}
\label{sec:Fragebogen_Konzept}

% TODO:
% - Von 96 auf 15 Fragen reduziert (DeVellis 2017)
% - Zwei Versionen: Vollversion (Brainstorming) + Kompakt (Workshop)
% - Tabelle mit Fragebogen-Reduktion
% - Begründung für Reduktion (Skalentwicklung nach DeVellis)


\subsection{Entity-Mapping-Workbook}
\label{sec:MappingWorkbook_Konzept}

% TODO:
% - JSON-Schema-Dokumentation für alle 6 Entity-Typen
% - Tabelle: Entity-Typen Übersicht
% - Für jede Entity: Felder, Datentypen, Required/Optional, Defaults
% - Beispiele für typische Konfigurationen


\subsection{Standard Operating Procedures (SOPs)}
\label{sec:SOPs_Konzept}

% TODO:
% - 3 SOPs beschreiben:
%   1. Kunden-Onboarding-Gespräch durchführen
%   2. Master Creator nutzen
%   3. JSONs in GRACE importieren
% - Format: Markdown, Git-versioniert
% - 6-Punkte-Validierungs-Checkliste


\subsection{Technische Legitimierung: Englischsprachiger Master Creator}
\label{sec:TechnischeLegitimierung}

% TODO:
% - LLM-Training primär auf englischen Daten
% - Ressourcenbeschränkungen (7B-Modell)
% - Performance-Unterschiede (30-40% besser bei Englisch)
% - Quellen: Lai et al. (2023), Ahuja et al. (2023), Zhang et al. (2024)

Die Wahl eines englischsprachigen Toolings für den Master Creator wird durch
mehrere technische und empirische Faktoren legitimiert. Large Language Models
werden primär auf englischsprachigen Korpora trainiert, was zu signifikanten
Leistungsunterschieden zwischen Englisch und anderen Sprachen führt. Studien
zeigen, dass kleinere Modelle (7B Parameter) bei nicht-englischen Sprachen
Leistungseinbußen von 30-40\% aufweisen \cite{llm:multilingual_performance2023,
llm:language_bias2023}. Diese Diskrepanz verstärkt sich bei Aufgaben der
strukturierten Ausgabegenerierung, wo Präzision und Formatkonformität
entscheidend sind \cite{llm:model_size_multilingual2024}.

Die Ressourcenbeschränkung auf ein 7B-Modell (Qwen2.5-Coder) resultiert aus
Datenhoheitsanforderungen: Lokales Deployment mit Ollama \cite{ollama:local_llm2024}
auf vorhandener Hardware limitiert die Modellgröße. Vergleichbare industrielle
Ansätze wie SmartAPS \cite{smartaps2025} nutzen ebenfalls leichtgewichtige
Modelle (Mistral-7B) und demonstrieren deren Anwendbarkeit für Operations
Management. Die Entscheidung für englischsprachige Prompts und JSON-Felder
folgt somit der technischen Notwendigkeit: Die Performance-Verbesserung bei
strukturierter Datengenerierung rechtfertigt den initalen Mehraufwand bei
der Prompt-Entwicklung und kompensiert die Ressourcenbeschränkung durch
optimale Modellauslastung.

